\documentclass[10pt,a4paper,ragged2e,withhyper]{altacv}

\geometry{left=1.25cm,right=1.25cm,top=1.5cm,bottom=1.5cm,columnsep=1.2cm}

\usepackage{paracol}
\usepackage[rm]{roboto}
\usepackage[defaultsans]{lato}
\usepackage[polish]{babel}

% Change the font if you want to, depending on whether
% you're using pdflatex or xelatex/lualatex
\renewcommand{\familydefault}{\sfdefault}

% Colors
\definecolor{SlateGrey}{HTML}{2E2E2E}
\definecolor{LightGrey}{HTML}{666666}
\definecolor{DarkPastelRed}{HTML}{450808}
\definecolor{PastelRed}{HTML}{8F0D0D}
\definecolor{GoldenEarth}{HTML}{E7D192}
\colorlet{name}{black}
\colorlet{tagline}{PastelRed}
\colorlet{heading}{DarkPastelRed}
\colorlet{headingrule}{GoldenEarth}
\colorlet{subheading}{PastelRed}
\colorlet{accent}{PastelRed}
\colorlet{emphasis}{SlateGrey}
\colorlet{body}{LightGrey}

% Fonts
\renewcommand{\namefont}{\Huge\rmfamily\bfseries}
\renewcommand{\personalinfofont}{\footnotesize}
\renewcommand{\cvsectionfont}{\LARGE\rmfamily\bfseries}
\renewcommand{\cvsubsectionfont}{\large\bfseries}
\renewcommand{\itemmarker}{{\small\textbullet}}
\renewcommand{\ratingmarker}{\faCircle}

\begin{document}

\name{Marek Maksimczyk}
\tagline{DevOps Engineer}
\personalinfo{%
  \email{marek.maksimczyk@mandos.net.pl}
  \homepage{mandos.net.pl}
  \github{mandos}
  \printinfo{\faLinkedin}{marek-maksimczyk}[https://linkedin.com/in/marek-maksimczyk-532a4358]
}

\makecvheader

\columnratio{0.6}
\begin{paracol}{2}

%% ====== LEWA KOLUMNA ======

\cvsection{Doświadczenie}

\cvevent{Niezależny Developer}{Open Source}{Listopad 2025 -- Obecnie}{}
\begin{itemize}
\item Zaprojektowanie i rozwój Open Transit Stack, open-source'owego monorepo do przetwarzania, walidacji i wizualizacji danych transportu publicznego (GTFS) z użyciem TypeScript, React i MapLibre GL.
\item Budowa infrastruktury AWS (S3, CloudFront, ACM, DNS) z wykorzystaniem Pulumi jako Infrastructure as Code.
\item Struktura projektu jako monorepo Nx ze ścisłym TypeScript, Vitest, testami E2E Playwright oraz ESLint/Prettier.
\end{itemize}

\divider

\cvevent{Senior DevOps Engineer}{Foundever}{Kwiecień 2024 -- Październik 2024}{}
\begin{itemize}
\item Zaprojektowanie samoobsługowego procesu provisioningu baz danych dla mikroserwisów, umożliwiającego zespołom deweloperskim tworzenie infrastruktury PostgreSQL bez udziału zespołu DevOps.
\item Przeprowadzenie migracji projektów GitLab na nową instancję, umożliwiając wyłączenie starszego serwera GitLab i standaryzację pipeline'ów CI/CD między zespołami.
\item Zarządzanie infrastrukturą AWS przy użyciu Terraform oraz usprawnienie procesów dostarczania, w tym pipeline'ów GitLab i konfiguracji ArgoCD/Helm.
\end{itemize}

\divider

\cvevent{Senior DevOps Engineer}{Cobiro}{Czerwiec 2021 -- Marzec 2024}{Kopenhaga, Dania}
\begin{itemize}
\item Kierowanie małym zespołem DevOps, nadzór nad infrastrukturą chmurową i procesami wdrożeń.
\item Osiągnięcie 50\% redukcji kosztów infrastruktury i operacyjnych poprzez optymalizację Kubernetes, dobór rozmiaru RDS oraz migrację z Datadog do Elastic Cloud.
\item Migracja infrastruktury AWS do modelu Infrastructure as Code (IaC) z wykorzystaniem Terraform, Terragrunt i Atlantis.
\item Administracja klastrami Kubernetes (EKS) z Helm, wdrożenie cluster autoscaler i KEDA do skalowania sterowanego zdarzeniami.
\item Budowa pipeline'ów CI/CD (CircleCI) i obrazów Docker dla zespołów frontend i backend.
\end{itemize}

\divider

\cvevent{Senior DevOps Engineer}{SeaChange Polska}{Styczeń 2019 -- Czerwiec 2021}{Warszawa, Polska}
\begin{itemize}
\item Koordynacja pracy małego zespołu DevOps oraz nadzór nad infrastrukturą i procesami automatyzacji.
\item Projektowanie i wdrożenie infrastruktury opartej na AWS z wykorzystaniem Terraform dla produktów firmy, obejmującej ECS, EKS, EC2, S3, CloudFront, KMS, IAM, ElastiCache (Redis) oraz Elasticsearch.
\item Tworzenie i konfiguracja systemów monitoringu (Prometheus, Grafana) oraz logowania (ELK, S3).
\item Opracowanie systemu dynamicznego tworzenia środowisk dla produktów firmy z wykorzystaniem Terraform i Jenkins.
\item Zarządzanie kontami AWS oraz konfiguracja infrastruktury sieciowej, w tym IAM i VPC, przy użyciu Terraform.
\end{itemize}

\divider

\cvevent{System Administrator}{Xstream A/S}{Styczeń 2015 -- Grudzień 2018}{Kopenhaga, Dania}
\begin{itemize}
\item Projektowanie i utrzymanie infrastruktury AWS z wykorzystaniem wewnętrznego systemu opartego na Python (Boto, Fabric), Puppet i Packer, a następnie migracja do CloudFormation.
\item Zarządzanie ponad 300 serwerami on-premise (zarówno fizycznymi, jak i wirtualnymi) przy użyciu Fabric i Puppet.
\item Zarządzanie i konfiguracja usług wewnętrznych, w tym: monitoringu (Nagios, Zabbix, Observium), repozytoriów kodu (SVN, Git, Gerrit), systemu logowania (ELK) oraz potoków CI/CD (Jenkins).
\end{itemize}

\divider

\cvevent{Senior PHP Developer}{Xstream A/S}{Styczeń 2012 -- Grudzień 2014}{Kopenhaga, Dania}
\begin{itemize}
\item Rozwój i utrzymanie systemów zarządzania treściami wideo oraz przetwarzania materiałów wideo z wykorzystaniem PHP 5 (Zend Framework, Symfony 2) i MySQL.
\end{itemize}

\divider

\cvevent{Wcześniejsze role: PHP Developer, System Administrator}{Różne}{2005 -- 2011}{}
\begin{itemize}
\item Budowa wielu aplikacji webowych i autorskich frameworków z użyciem PHP 5, MySQL oraz różnych frameworków (Kohana, Yii, Zend).
\item Implementacja rozwiązań do przetwarzania i wyszukiwania danych z użyciem Lucene i Solr.
\item Wsparcie IT, w tym zarządzanie siecią i administracja serwerami Windows.
\end{itemize}

%% ====== PRAWA KOLUMNA ======

\switchcolumn

\cvsection{Profil}
Senior DevOps Engineer z ponad 10-letnim doświadczeniem w infrastrukturze chmurowej i zapleczem programistycznym. Doświadczony w kierowaniu małymi zespołami DevOps, migracji infrastruktury do IaC oraz optymalizacji kosztów AWS. Specjalizacja w Terraform, Kubernetes i automatyzacji CI/CD.

\cvsection{Umiejętności techniczne}

\cvtag{AWS}
\cvtag{GCP}\\
\cvtag{Terraform}
\cvtag{Terragrunt}
\cvtag{Pulumi}\\
\cvtag{Docker}
\cvtag{Kubernetes}
\cvtag{Helm}\\
\cvtag{GitLab CI}
\cvtag{CircleCI}
\cvtag{ArgoCD}\\
\cvtag{Datadog}
\cvtag{Prometheus}
\cvtag{Grafana}
\cvtag{ELK Stack}\\
\cvtag{Python}
\cvtag{Bash}
\cvtag{TypeScript}

\cvsection{Języki}
\begin{itemize}
\item Polski -- Ojczysty
\item Angielski -- C1 (Biegłość zawodowa)
\item Japoński -- A1 (Początkujący)
\end{itemize}

\cvsection{Edukacja}
\cvevent{Politechnika Łódzka}{Wydział Fizyki Technicznej, Informatyki i Matematyki Stosowanej}{2001 -- 2006}{Łódź, Polska}

\end{paracol}

\end{document}
